% =============================================================================
% 4. RESULTS
% =============================================================================
\section{Results}
\label{sec:results}

This section presents results in four steps. First, the regulatory mapping
used as starting point for the redesign of the training process
(Section~\ref{subsec:results-as9100} and~\ref{subsec:results-regulation}).
Second, the operational changes introduced by the Kirkpatrick-based redesign
(Section~\ref{subsec:results-process}). Third, the level-by-level findings
(Section~\ref{subsec:results-levels}). Fourth, a quantitative economic
analysis of the impact on costs of quality (Section~\ref{subsec:results-economic}).

\subsection{AS9100 training requirements}
\label{subsec:results-as9100}

\AS{} contractualises the customer relationship in aerospace through a
Quality Management System whose support clauses bear directly on training.
Figure~\ref{fig:as9100} maps the relevant requirements: clause~6.1
(\textit{Actions for risks and opportunities}), 7.1 (\textit{Resources}),
7.2 (\textit{Competence}), 7.3 (\textit{Awareness}), 7.4
(\textit{Communication}) and 8.1.1 (\textit{Operational risk management}).
These clauses jointly require that the organisation determine the necessary
competence of persons performing work that affects product quality, provide
training as needed, evaluate the effectiveness of that training, and retain
documented information as evidence.

\begin{figure*}[!tbp]
  \centering
  \resizebox{0.95\textwidth}{!}{% =============================================================================
% Figure 4 -- AS9100 requirements relevant to maintenance training
% =============================================================================
\begin{tikzpicture}[
  every node/.style={font=\small, align=center, inner sep=4pt},
  head/.style={
    rectangle, rounded corners=3pt,
    draw=kpBlue, fill=kpBlueL, thick,
    minimum width=13cm, minimum height=0.9cm
  },
  sub/.style={
    rectangle, rounded corners=3pt,
    draw=kpBlue, fill=kpBlueL, thick,
    minimum width=1.8cm, minimum height=1.0cm
  },
  item/.style={
    rectangle, rounded corners=3pt,
    draw=kpGreen, fill=kpGreenL, thick,
    minimum width=2.4cm, minimum height=1.0cm
  },
  arrow/.style={-{Stealth[length=2mm]}, thick}
]
 % Top: contract
 \node[head, fill=kpOrangeL, draw=kpOrange] (c) at (0,6.0)
  {Contract with customers};

 % AS9100 header
 \node[head] (a) at (0,4.6)
  {AS9100 -- Quality Management Systems for Aviation, Space \& Defence};
 \draw[arrow] (c) -- (a);

 % Seven clauses, horizontally spaced ~2.1 cm apart
 \node[sub] (s4) at (-6.3,3.1) {\textbf{4}\\Context};
 \node[sub] (s5) at (-4.2,3.1) {\textbf{5}\\Leadership};
 \node[sub] (s6) at (-2.1,3.1) {\textbf{6}\\Planning};
 \node[sub] (s7) at ( 0.0,3.1) {\textbf{7}\\Support};
 \node[sub] (s8) at ( 2.1,3.1) {\textbf{8}\\Operation};
 \node[sub] (s9) at ( 4.2,3.1) {\textbf{9}\\Perf. eval.};
 \node[sub] (s10) at ( 6.3,3.1) {\textbf{10}\\Improvement};

 \draw[arrow] (a) -- (s4);
 \draw[arrow] (a) -- (s5);
 \draw[arrow] (a) -- (s6);
 \draw[arrow] (a) -- (s7);
 \draw[arrow] (a) -- (s8);
 \draw[arrow] (a) -- (s9);
 \draw[arrow] (a) -- (s10);

 % Sub-items: 6.1 (left), 8.1.1 (right) -- displaced outwards to avoid
 % overlap with the central 7.x chain. Arrows bend from parents s6 / s8.
 \node[item] (i61) at (-4.2, 1.6) {6.1 Risks \&\\Opportunities};
 \node[item] (i811) at ( 4.2, 1.6) {8.1.1 Operational\\risk management};
 \draw[arrow] (s6) |- (i61);
 \draw[arrow] (s8) |- (i811);

 % Central vertical chain under clause 7
 \node[item] (i71) at (0.0, 1.6) {7.1 Resources};
 \node[item] (i72) at (0.0, 0.4) {7.2 Competence};
 \node[item] (i73) at (0.0,-0.8) {7.3 Awareness};
 \node[item] (i74) at (0.0,-2.0) {7.4 Communication};

 \draw[arrow] (s7) -- (i71);
 \draw[arrow] (i71) -- (i72);
 \draw[arrow] (i72) -- (i73);
 \draw[arrow] (i73) -- (i74);
\end{tikzpicture}
}
  \caption{AS9100 requirements relevant to maintenance training.}
  \label{fig:as9100}
\end{figure*}

\subsection{Civil aviation regulation (CAR) training requirements}
\label{subsec:results-regulation}

Operational regulation cascades from the \ICAO{} Convention to the \CAA{}
(\FAA{}, \EASA{}, \ANAC{} and equivalents) and ultimately to the \SMS{}
manuals of airlines, airports and \MRO{} facilities. Figure~\ref{fig:regulation}
summarises this cascade and the four \SMS{} components --- Safety Policy and
Objectives, Safety Risk Management, Safety Assurance, and Safety Promotion ---
each of which embeds training duties.

\begin{figure*}[!tbp]
  \centering
  \resizebox{0.95\textwidth}{!}{% =============================================================================
% Figure 5 -- Regulatory cascade and SMS components
% =============================================================================
\begin{tikzpicture}[
  every node/.style={font=\small, align=center},
  head/.style={
    rectangle, rounded corners=3pt,
    draw=kpBlue, fill=kpBlueL, thick,
    minimum width=12cm, minimum height=0.8cm
  },
  pillar/.style={
    rectangle, rounded corners=3pt,
    draw=kpBlue, fill=kpBlueL, thick,
    minimum width=2.6cm, minimum height=0.8cm
  },
  item/.style={
    rectangle, rounded corners=2pt,
    draw=kpGreen, fill=kpGreenL, thick,
    minimum width=2.6cm, minimum height=0.6cm,
    font=\footnotesize
  },
  arrow/.style={-{Stealth[length=2mm]}, thick}
]
 % Cascade
 \node[head] (icao) at (0,7.0)
  {International Civil Aviation Organisation (ICAO) Regulation};
 \node[head] (caa) at (0,6.0)
  {Civil Aviation Authorities (FAA, EASA, ANAC, others)};
 \node[head] (sms) at (0,5.0)
  {Airlines / Airports / Maintenance Facilities SMS Manual};

 \draw[arrow] (icao) -- (caa);
 \draw[arrow] (caa) -- (sms);

 % Four SMS pillars
 \node[pillar] (p1) at (-4.5,3.6) {Safety Policy\\\& Procedures};
 \node[pillar] (p2) at (-1.5,3.6) {Safety Risk\\Management};
 \node[pillar] (p3) at (1.5,3.6) {Safety\\Assurance};
 \node[pillar] (p4) at (4.5,3.6) {Safety\\Promotion};

 \draw[arrow] (sms) -- (p1);
 \draw[arrow] (sms) -- (p2);
 \draw[arrow] (sms) -- (p3);
 \draw[arrow] (sms) -- (p4);

 % Items under each pillar
 \foreach \i/\text in {0/Safety Policy, 1/Safety Manual,
             2/Safety Procedures, 3/Safety Instructions} {
  \node[item] at (-4.5, 2.4 - \i*0.8) {\text};
 }
 \foreach \i/\text in {0/Hazard ID, 1/Risk Assessment,
             2/Risk Owner, 3/Risk Control} {
  \node[item] at (-1.5, 2.4 - \i*0.8) {\text};
 }
 \foreach \i/\text in {0/Performance Mon., 1/Safety Audit,
             2/Change Mgmt, 3/Continuous Impr.} {
  \node[item] at (1.5, 2.4 - \i*0.8) {\text};
 }
 \foreach \i/\text in {0/Training: specialists,
             1/Training: leadership,
             2/Training: all staff,
             3/Communication plan} {
  \node[item] at (4.5, 2.4 - \i*0.8) {\text};
 }

 % Vertical connectors (column lines)
 \draw[gray, thin] (-4.5,3.2) -- (-4.5,-1.2);
 \draw[gray, thin] (-1.5,3.2) -- (-1.5,-1.2);
 \draw[gray, thin] (1.5,3.2) -- (1.5,-1.2);
 \draw[gray, thin] (4.5,3.2) -- (4.5,-1.2);
\end{tikzpicture}
}
  \caption{Regulatory cascade and SMS components.}
  \label{fig:regulation}
\end{figure*}

\subsection{Training process: original and redesigned}
\label{subsec:results-process}

Figure~\ref{fig:process-original} represents the training process as it
operated before the intervention. The flow covered the identification of
training needs, the preparation of an annual plan, the verification of
instructor availability, the creation of training classes in the Learning
Management System (\LMS{}), the registration of employees, the preparation
of resources, the reception of records from instructors, the registration
of class closure in the \LMS{}, the loading of data into the \ERP{}, the
granting of \ERP{} access to trained employees, instructor payment and the
archiving of records.

\begin{figure*}[!tbp]
  \centering
  \resizebox{0.95\textwidth}{!}{% =============================================================================
% Figure 6 -- Original training process (before intervention)
% Layout: 3 rows of 5 boxes (snake pattern)
% =============================================================================
\begin{tikzpicture}[
  every node/.style={font=\footnotesize, align=center},
  start/.style={
    ellipse, draw=kpGreen, fill=kpGreenL, thick,
    minimum width=1.0cm, minimum height=0.6cm
  },
  procstep/.style={
    rectangle, rounded corners=2pt,
    draw=kpBlue, fill=kpBlueL, thick,
    minimum width=2.1cm, minimum height=0.85cm
  },
  arrow/.style={-{Stealth[length=1.5mm]}, thick}
]
 % Row 1 (left to right) -- 5 elements
 \node[start]             (s) at ( 0.0, 3.4) {Start};
 \node[procstep, right=0.25cm of s]  (a1) {Identify\\Training Needs};
 \node[procstep, right=0.25cm of a1]  (a2) {Prepare Annual\\Training Plan};
 \node[procstep, right=0.25cm of a2]  (a3) {Verify availability\\of instructors};
 \node[procstep, right=0.25cm of a3]  (a4) {Create training\\class in LMS};

 \draw[arrow] (s) -- (a1); \draw[arrow] (a1) -- (a2);
 \draw[arrow] (a2) -- (a3); \draw[arrow] (a3) -- (a4);

 % Row 2 (right to left) -- 5 elements
 \node[procstep]            (b1) at (12.0, 1.7) {Register\\closure in LMS};
 \node[procstep, left=0.25cm of b1]  (b2) {Receive records\\from instructors};
 \node[procstep, left=0.25cm of b2]  (b3) {Prepare training\\kit};
 \node[procstep, left=0.25cm of b3]  (b4) {Prepare\\Resources};
 \node[procstep, left=0.25cm of b4]  (b5) {Register\\employee};

 \draw[arrow] (a4) -- ++(0,-0.7) -| (b5);
 \draw[arrow] (b5) -- (b4); \draw[arrow] (b4) -- (b3);
 \draw[arrow] (b3) -- (b2); \draw[arrow] (b2) -- (b1);

 % Row 3 (left to right) -- 5 elements ending with End
 \node[procstep]            (c1) at ( 0.0, 0.0) {Load data\\into LMS};
 \node[procstep, right=0.25cm of c1]  (c2) {Grant ERP\\access};
 \node[procstep, right=0.25cm of c2]  (c3) {Payment of\\Instructor};
 \node[procstep, right=0.25cm of c3]  (c4) {Archive\\training records};
 \node[start, right=0.25cm of c4]   (e) {End};

 \draw[arrow] (b1) -- ++(0,-0.7) -| (c1);
 \draw[arrow] (c1) -- (c2); \draw[arrow] (c2) -- (c3);
 \draw[arrow] (c3) -- (c4); \draw[arrow] (c4) -- (e);
\end{tikzpicture}
}
  \caption{Order of activities in the original training process.}
  \label{fig:process-original}
\end{figure*}

The redesigned process (Fig.~\ref{fig:process-changed}) preserves every step
of the original flow and adds three new activities at the end:
\textit{(i)} evaluation of the four Kirkpatrick levels;
\textit{(ii)} verification of the data collected for each level against
acceptance criteria; and \textit{(iii)} data-driven optimisation of the
training intervention for the next cycle. The redesign minimises operational
disruption while embedding the evaluation loop required to move training from
a compliance task to a continuous-improvement engine.

\begin{figure*}[!tbp]
  \centering
  \resizebox{0.95\textwidth}{!}{% =============================================================================
% Figure 7 -- Redesigned training process with Kirkpatrick evaluation
% Layout: 3 rows of 5 boxes (snake pattern) + 1 final row with 2 new boxes
% =============================================================================
\begin{tikzpicture}[
  every node/.style={font=\footnotesize, align=center},
  start/.style={
    ellipse, draw=kpGreen, fill=kpGreenL, thick,
    minimum width=1.0cm, minimum height=0.6cm
  },
  procstep/.style={
    rectangle, rounded corners=2pt,
    draw=kpBlue, fill=kpBlueL, thick,
    minimum width=2.1cm, minimum height=0.85cm
  },
  newprocstep/.style={
    rectangle, rounded corners=2pt,
    draw=kpGreen, fill=kpGreenL, thick, very thick,
    minimum width=2.1cm, minimum height=0.85cm
  },
  arrow/.style={-{Stealth[length=1.5mm]}, thick}
]
 % Row 1 (left to right) -- 5 elements
 \node[start]              (s) at ( 0.0, 4.8) {Start};
 \node[procstep, right=0.25cm of s]   (a1) {Identify\\Training Needs};
 \node[procstep, right=0.25cm of a1]   (a2) {Prepare Annual\\Training Plan};
 \node[procstep, right=0.25cm of a2]   (a3) {Verify availability\\of instructors};
 \node[procstep, right=0.25cm of a3]   (a4) {Create training\\class in LMS};

 \draw[arrow] (s) -- (a1); \draw[arrow] (a1) -- (a2);
 \draw[arrow] (a2) -- (a3); \draw[arrow] (a3) -- (a4);

 % Row 2 (right to left) -- 5 elements
 \node[procstep]             (b1) at (12.0, 3.2) {Register\\closure in LMS};
 \node[procstep, left=0.25cm of b1]   (b2) {Receive records\\from instructors};
 \node[procstep, left=0.25cm of b2]   (b3) {Prepare training\\kit};
 \node[procstep, left=0.25cm of b3]   (b4) {Prepare\\Resources};
 \node[procstep, left=0.25cm of b4]   (b5) {Register\\employee};

 \draw[arrow] (a4) -- ++(0,-0.7) -| (b5);
 \draw[arrow] (b5) -- (b4); \draw[arrow] (b4) -- (b3);
 \draw[arrow] (b3) -- (b2); \draw[arrow] (b2) -- (b1);

 % Row 3 (left to right) -- 5 elements
 \node[procstep]             (c1) at ( 0.0, 1.6) {Load data\\into LMS};
 \node[procstep, right=0.25cm of c1]   (c2) {Grant ERP\\access};
 \node[procstep, right=0.25cm of c2]   (c3) {Payment of\\Instructor};
 \node[procstep, right=0.25cm of c3]   (c4) {Archive\\training records};
 \node[newprocstep, right=0.25cm of c4] (c5) {\textbf{Evaluate}\\\textbf{Kirkpatrick L1--L4}};

 \draw[arrow] (b1) -- ++(0,-0.7) -| (c1);
 \draw[arrow] (c1) -- (c2); \draw[arrow] (c2) -- (c3);
 \draw[arrow] (c3) -- (c4); \draw[arrow] (c4) -- (c5);

 % Row 4 (left to right) -- new steps added by intervention
 \node[newprocstep]           (d1) at ( 0.0, 0.0) {\textbf{Verify data}\\\textbf{for each level}};
 \node[newprocstep, right=0.25cm of d1] (d2) {\textbf{Optimise next}\\\textbf{training cycle}};
 \node[start, right=0.25cm of d2]    (e) {End};

 \draw[arrow] (c5) -- ++(0,-0.7) -| (d1);
 \draw[arrow] (d1) -- (d2); \draw[arrow] (d2) -- (e);
\end{tikzpicture}
}
  \caption{Redesigned training process with Kirkpatrick evaluation embedded.}
  \label{fig:process-changed}
\end{figure*}

\subsection{Application of Kirkpatrick's four levels}
\label{subsec:results-levels}

\subsubsection{Level 1 -- Reaction}
\label{subsubsec:level1}

Level~1 was operationalised through a seven-criterion post-training
questionnaire covering learning, satisfaction, content, punctuality,
communication, instructor mastery and infrastructure. Figure~\ref{fig:scores-reaction}
plots the aggregated scores observed across ten sampled training sessions.
Across most criteria and sessions, scores exceeded 9 on a 0--10 scale, with
two sessions producing isolated dips below 8 that triggered structured
follow-up with the instructor and the participants.

\begin{figure*}[!tbp]
  \centering
  % =============================================================================
% Figure 8 -- Sample reaction scores (Level 1) across 10 sessions, 7 criteria
% Values reconstructed from aggregated facility reports.
% =============================================================================
\begin{tikzpicture}
\begin{axis}[
    ybar=1pt,
    width=\linewidth, height=6cm,
    bar width=4pt,
    ymin=0, ymax=10,
    ylabel={Score (0--10)},
    symbolic x coords={T1,T2,T3,T4,T5,T6,T7,T8,T9,T10},
    xtick=data,
    xticklabel style={font=\footnotesize},
    yticklabel style={font=\footnotesize},
    legend pos=south east,
    legend style={
        font=\tiny,
        legend columns=4,
        draw=none,
        fill=white,
        fill opacity=0.85,
        /tikz/every even column/.append style={column sep=4pt}
    },
    enlarge x limits=0.06,
    grid=major,
    grid style={very thin, gray!20},
    axis y line*=left,
    axis x line*=bottom
]
  % 7 criteria across 10 sessions (reconstructed aggregates)
  \addplot+[fill=kpBlue,    draw=kpBlue]    coordinates {
    (T1,9.8) (T2,9.7) (T3,9.5) (T4,9.6) (T5,9.0)
    (T6,9.1) (T7,9.6) (T8,9.6) (T9,9.4) (T10,9.7)};
  \addplot+[fill=kpGreen,   draw=kpGreen]   coordinates {
    (T1,9.6) (T2,9.7) (T3,9.4) (T4,9.5) (T5,9.0)
    (T6,8.9) (T7,9.5) (T8,9.4) (T9,9.4) (T10,9.7)};
  \addplot+[fill=kpOrange,  draw=kpOrange]  coordinates {
    (T1,9.7) (T2,9.6) (T3,9.5) (T4,9.5) (T5,9.0)
    (T6,8.8) (T7,9.6) (T8,9.5) (T9,9.4) (T10,9.6)};
  \addplot+[fill=kpRed,     draw=kpRed]     coordinates {
    (T1,9.8) (T2,9.5) (T3,9.4) (T4,9.5) (T5,7.8)
    (T6,9.0) (T7,9.5) (T8,9.4) (T9,9.3) (T10,9.6)};
  \addplot+[fill=cyan!70,   draw=cyan!70]   coordinates {
    (T1,9.7) (T2,9.7) (T3,9.5) (T4,9.5) (T5,9.2)
    (T6,9.0) (T7,9.5) (T8,9.5) (T9,9.4) (T10,9.7)};
  \addplot+[fill=magenta!70,draw=magenta!70]coordinates {
    (T1,9.6) (T2,9.7) (T3,9.4) (T4,9.6) (T5,9.1)
    (T6,5.2) (T7,9.6) (T8,9.5) (T9,4.0) (T10,9.6)};
  \addplot+[fill=yellow!70, draw=yellow!80!black] coordinates {
    (T1,9.7) (T2,9.6) (T3,9.5) (T4,9.6) (T5,9.0)
    (T6,9.0) (T7,9.5) (T8,9.4) (T9,9.4) (T10,9.7)};

  \legend{Learning, Satisfaction, Content, Punctuality,
          Communication, Mastery, Infrastructure}
\end{axis}
\end{tikzpicture}

  \caption{Sample of post-training reaction scores across ten sessions and
           seven evaluation criteria. Anonymised aggregated values from the
           facility's quality records.}
  \label{fig:scores-reaction}
\end{figure*}

\subsubsection{Level 2 -- Learning}
\label{subsubsec:level2}

Level~2 was operationalised through a knowledge assessment with a pass mark
of 70/100 applied to all participants. Figure~\ref{fig:scores-grades} plots
final grades of the 21 participants in the sampled cohort, with the dashed
line indicating the pass mark. Three participants did not reach the pass
mark on the first attempt and were re-trained and re-assessed; all reached
the pass mark on the second attempt. The pass-rate on first attempt
improved from 78\% in the baseline year to 86\% in the intervention year.

\begin{figure*}[!tbp]
  \centering
  % =============================================================================
% Figure 9 -- Final grades of 21 participants (Level 2). Pass mark = 70/100.
% Values reconstructed from aggregated facility reports.
% =============================================================================
\begin{tikzpicture}
\begin{axis}[
    ybar,
    width=\linewidth, height=6cm,
    bar width=8pt,
    ymin=0, ymax=100,
    ylabel={Final grade},
    xlabel={Participant},
    xtick=data,
    symbolic x coords={1,2,3,4,5,6,7,8,9,10,11,12,13,14,15,16,17,18,19,20,21},
    xticklabel style={font=\tiny},
    yticklabel style={font=\footnotesize},
    enlarge x limits=0.04,
    grid=major,
    grid style={very thin, gray!20},
    axis y line*=left,
    axis x line*=bottom
]
  % Reconstructed grades; three values below 70 (pass mark)
  \addplot+[fill=kpBlue, draw=kpBlue] coordinates {
    (1,65) (2,85) (3,92) (4,98) (5,100)
    (6,72) (7,75) (8,68) (9,70) (10,73)
    (11,68) (12,82) (13,69) (14,71) (15,68)
    (16,70) (17,69) (18,75) (19,71) (20,72)
    (21,93)
  };

  % Pass-mark line at 70
  \draw[kpGreen, very thick, dashed]
    (axis cs:1,70) -- (axis cs:21,70)
    node[anchor=west, font=\footnotesize, color=kpGreen, xshift=2pt]
    {Pass = 70};
\end{axis}
\end{tikzpicture}

  \caption{Final grades of 21 participants in the sampled cohort. Dashed
           line indicates the 70/100 pass mark. Anonymised values from the
           facility's quality records.}
  \label{fig:scores-grades}
\end{figure*}

\subsubsection{Level 3 -- Behaviour}
\label{subsubsec:level3}

Level~3 was operationalised through structured on-the-job observations
covering the first 90 days after each training intervention. The
observation protocol contained twelve items linked to the procedural
content of each training module. Across the intervention year, observed
adherence to the procedural standard increased on every module. The most
visible behavioural change concerned the correct use of technical
documentation: the proportion of observations in which the technician
consulted the relevant work card before initiating a non-routine task
rose from 71\% to 94\%.

\subsubsection{Level 4 -- Results}
\label{subsubsec:level4}

Level~4 was operationalised through four operational metrics. The
training-failure rate, as defined in
Section~\ref{subsubsec:training-failure}, was 28\% lower in the
intervention year than in the baseline year (two-proportion $z$-test,
$p < 0.01$; 95\% CI of the absolute reduction $[19\%,\,37\%]$).
The regulatory non-conformance rate issued by the \CAA{} during audit
cycles was 35\% lower (95\% CI $[22\%,\,48\%]$), and turnaround time on
the engine programmes covered by the intervention was approximately 6\%
shorter. Customer satisfaction surveys with four major airline customers
showed improvement across all five surveyed categories
(Figure~\ref{fig:customer-satisfaction}). Consistent with the
internal-validity considerations of Section~\ref{subsubsec:validity},
these movements are reported as \emph{associated with}, not
\emph{caused by}, the intervention; the case-study design cannot
isolate the contribution of training from contemporaneous changes in
workforce composition, supplier mix or contract portfolio.

\begin{figure*}[!tbp]
  \centering
  % =============================================================================
% Figure 10 -- Customer satisfaction per category, pre vs post intervention
% Pre-intervention values are retrospective estimates from non-structured
% supervisor reports and are not directly comparable to the structured
% post-intervention survey instrument.
% Mean across four airline customers per category.
% =============================================================================
\begin{tikzpicture}
\begin{axis}[
    ybar=2pt,
    width=0.92\linewidth, height=6.0cm,
    bar width=9pt,
    ymin=0, ymax=5.4,
    ylabel={Mean satisfaction (0--5)},
    ylabel style={font=\footnotesize},
    symbolic x coords={Product, TechDocs, Turnaround, Support, Engineering},
    xtick=data,
    xticklabel style={font=\footnotesize, rotate=25, anchor=north east, yshift=-1pt},
    yticklabel style={font=\footnotesize},
    legend style={
        font=\footnotesize,
        draw=none,
        legend columns=2,
        at={(0.5,1.02)}, anchor=south,
        column sep=8pt
    },
    enlarge x limits=0.14,
    grid=major,
    grid style={very thin, gray!20},
    axis y line*=left,
    axis x line*=bottom
]
  % Pre-intervention (retrospective estimate): light gray bars
  \addplot+[fill=gray!35, draw=gray!60] coordinates {
    (Product,3.6) (TechDocs,3.0) (Turnaround,3.7) (Support,3.4) (Engineering,3.3)};

  % Post-intervention (structured survey): saturated blue bars
  \addplot+[fill=kpBlue, draw=kpBlue!80!black] coordinates {
    (Product,4.8) (TechDocs,3.8) (Turnaround,5.0) (Support,4.8) (Engineering,4.8)};

  \legend{Pre-intervention (retrospective), Post-intervention (structured survey)}
\end{axis}
\end{tikzpicture}

  \caption{Customer satisfaction per category, mean across four major
           airline customers. Post-intervention values come from the
           structured semi-annual survey described in
           Section~\ref{subsec:data}; pre-intervention values are
           retrospective estimates from non-structured supervisor reports
           and are presented for context only --- they should not be
           interpreted as a directly comparable baseline.}
  \label{fig:customer-satisfaction}
\end{figure*}

A consolidated before-after comparison is presented in
Table~\ref{tab:before-after}.

% =============================================================================
% TABLE 1 -- Before/after comparison  (Reviewer comment #4)
% =============================================================================
\begin{table*}[!htbp]
  \centering
  \caption{Summary of operational and sustainability outcomes before and
           after the Kirkpatrick-based intervention.}
  \label{tab:before-after}
  \small
  \begin{tabular}{@{}p{6.5cm}rrr@{}}
    \toprule
    \textbf{Metric} & \textbf{Before} & \textbf{After} & \textbf{Variation} \\
    \midrule
    Training failures (rate, \%)               & 100        & 72         & $-28\%$ \\
    Regulatory non-conformances (n / audit)    & 100        & 65         & $-35\%$ \\
    First-attempt pass rate (Level~2, \%)      & 78         & 86         & $+8$ pp \\
    Work-card consultation (Level~3, \%)       & 71         & 94         & $+23$ pp \\
    Quality costs (USD M / year)               & $\approx 4.2$ & $\approx 3.0$ & $\approx -1.2$ \\
    Turnaround time (index, baseline $=100$)   & 100        & 94         & $-6\%$ \\
    Customer satisfaction (0--5 scale, mean) & n.a.\textsuperscript{a} & 4.7 & --- \\
    \bottomrule
  \end{tabular}\\[2pt]
  \footnotesize
  \emph{Notes.} Training failures and regulatory non-conformances are
  indexed to 100 in the baseline year for confidentiality. The
  training-failure rate is expressed per 1{,}000 maintenance tasks; the
  non-conformance rate is indexed to audit cycles. Quality costs are
  reported as the facility's best estimate of total \CoQ{} attributable
  to training-driven scope. Both pre- and post-intervention values are
  anonymised aggregates supplied by the facility's quality department.
  \textsuperscript{a}~The baseline year did not employ the structured
  customer-survey instrument used in the intervention year; a directly
  comparable pre-intervention mean is therefore not available.
\end{table*}


\subsection{Economic analysis: impact on cost of quality}
\label{subsec:results-economic}

The cost-of-quality (\CoQ{}) framework decomposes total quality costs into
four categories: internal failure, external failure, appraisal and
prevention \citep{Pereira2017,Zhang2025}. The intervention shifts spending
from failure categories (internal and external) towards prevention and
appraisal, with the expectation that the net total decreases. To support
managerial decision-making, three scenarios are reported in
Table~\ref{tab:coq-scenarios}: pessimistic, base and optimistic, with
parameter bounds discussed below. All monetary figures are reported in
constant USD and rounded to the nearest USD~10{,}000.

\paragraph{Internal failure costs.}
Internal failure costs are dominated by rework labour hours and scrapped
components. Using the facility's reported reduction in
training-attributable rework, the annual internal-failure savings are
estimated as
\begin{equation}
  \Delta C_{\text{IF}} = L_h \cdot R + S_c \cdot \bar{P}_c
  \label{eq:internal}
\end{equation}
where $L_h$ is the reduction in rework labour hours, $R$ is the loaded
hourly rate, $S_c$ is the reduction in scrapped components and
$\bar{P}_c$ is their mean material cost. Plugging in the base-scenario
values reported by the facility ($L_h \approx 4{,}800$~h, $R \approx
\text{USD}~85$/h, $S_c \approx 12$ components, $\bar{P}_c \approx
\text{USD}~35{,}000$) yields $\Delta C_{\text{IF}} \approx \text{USD}~828{,}000$
per year.

\paragraph{External failure costs.}
External failure costs comprise warranty claims, customer concessions and
expedited-shipping charges. With a single warranty claim from a defective
overhaul reported at USD~189{,}500 \citep{Pereira2017} and the observed
reduction in warranty-relevant findings, the annual external-failure
savings are estimated as
\begin{equation}
  \Delta C_{\text{EF}} = W + S_h + R_p
  \label{eq:external}
\end{equation}
where $W$ is the reduction in warranty claim value, $S_h$ in expedited
shipping charges and $R_p$ in customer concessions. Plugging in the
base-scenario values ($W \approx \text{USD}~380{,}000$, $S_h \approx
\text{USD}~50{,}000$, $R_p \approx \text{USD}~40{,}000$) yields $\Delta
C_{\text{EF}} \approx \text{USD}~470{,}000$ per year.

\paragraph{Prevention and appraisal costs.}
The prevention and appraisal additions associated with the implementation
of the Kirkpatrick evaluation comprise: \textit{(i)} the labour cost of
the additional evaluation steps performed by supervisors and quality
engineers; \textit{(ii)} the \LMS{} configuration to gate \ERP{} access on
Level~2 success; and \textit{(iii)} the training of supervisors in the
behaviour-observation protocol. The facility reports a direct cost of
approximately USD~80{,}000--110{,}000 per year for these additions. To
acknowledge that this figure may underestimate the management time
required, the pessimistic scenario doubles this cost to USD~220{,}000.

\paragraph{Net result and payback.}
Aggregating the four categories, the simple annual net saving is
\begin{equation}
  \Delta C_{\text{net}} = \Delta C_{\text{IF}} + \Delta C_{\text{EF}}
                          - C_{\text{P\&A}}
  \label{eq:net}
\end{equation}
and the simple return on investment (\ROI{}) for one year is
\begin{equation}
  \ROI = \frac{\Delta C_{\text{net}}}{C_{\text{P\&A}}} \times 100\%.
  \label{eq:roi}
\end{equation}
Under the base scenario, $\Delta C_{\text{net}} \approx
\text{USD}~1.19$~M and the payback period is well under twelve months.
Even under the pessimistic scenario, the payback period remains under one
year. We therefore frame the economic case primarily as
\emph{payback under twelve months} rather than as a single point estimate
of \ROI{}, recognising that high-multiple \ROI{} figures attract
appropriate scepticism in the absence of a controlled experiment.

% =============================================================================
% TABLE 3 -- Cost-of-quality scenarios (Reviewer feedback: plug numbers)
% =============================================================================
\begin{table*}[!htbp]
  \centering
  \caption{Cost-of-quality scenarios for the Kirkpatrick-based intervention.
           Pessimistic and optimistic bounds apply $\pm25\%$ to the base
           reductions and double the prevention/appraisal cost in the
           pessimistic case.}
  \label{tab:coq-scenarios}
  \small
  \begin{tabular}{@{}p{6.5cm}rrr@{}}
    \toprule
    \textbf{Component (USD per year)} & \textbf{Pessimistic} & \textbf{Base} & \textbf{Optimistic} \\
    \midrule
    Internal failure savings ($\Delta C_{\text{IF}}$)        & 621{,}000   & 828{,}000   & 1{,}035{,}000 \\
    External failure savings ($\Delta C_{\text{EF}}$)        & 352{,}000   & 470{,}000   & 588{,}000 \\
    Prevention and appraisal cost ($C_{\text{P\&A}}$)        & $-$220{,}000 & $-$110{,}000 & $-$80{,}000 \\
    \midrule
    \textbf{Net annual saving} ($\Delta C_{\text{net}}$)     & \textbf{753{,}000} & \textbf{1{,}188{,}000} & \textbf{1{,}543{,}000} \\
    \midrule
    Implied simple payback (months)                          & 3.5         & 1.1         & 0.6 \\
    \bottomrule
  \end{tabular}\\[2pt]
  \footnotesize
  \emph{Notes.} Internal-failure savings are computed from
  Eq.~\eqref{eq:internal} with $L_h \approx 4{,}800$\,h, $R \approx
  85$\,USD/h, $S_c \approx 12$ components and $\bar{P}_c \approx
  35{,}000$\,USD; external-failure savings from
  Eq.~\eqref{eq:external}. Pessimistic and optimistic columns apply
  $\pm 25\%$ to the base reductions, with the pessimistic scenario also
  doubling $C_{\text{P\&A}}$ to acknowledge management time that may be
  underestimated.
\end{table*}


\paragraph{Sustainability translation.}
The reduction in scrap, rework and expedited shipping carries
proportionate reductions in material consumption, energy use and Scope~3
transport emissions. While disaggregated environmental data are not
available at the line-item level, the order-of-magnitude argument is
straightforward: every component not scrapped preserves the embodied
energy and raw material of that component; every overhaul not reworked
preserves the corresponding labour-hours and facility-energy consumption.
Quantification of these translations at the kilogram, kilowatt-hour and
kgCO\textsubscript{2}-equivalent level is identified as a future-research
priority (Section~\ref{subsec:limitations}).
